\documentclass[12pt,a4paper]{article}
\usepackage[utf8]{vietnam}
\usepackage{amsmath,amssymb}
\usepackage{geometry}
\geometry{
  top=12mm,
  bottom=12mm,
  left=14mm,
  right=14mm,
  footskip=6mm,
  headheight=14pt
}

\usepackage{tikz}
\usetikzlibrary{arrows.meta,calc,shapes.geometric}
\usepackage[most]{tcolorbox}
\usepackage{fontawesome}
\usepackage{fancyhdr}
\usepackage{enumitem}

\pagestyle{fancy}
\fancyhf{}
\fancyhead[L]{\footnotesize\color{teal!80!black}\faBook\ \textbf{SỔ TAY KIẾN THỨC TOÁN 3}}
\fancyhead[R]{\footnotesize\color{teal!80!black}\textbf{Lớp Toán Thầy Lĩnh -- loptoanthaylinh.vn}}
\fancyfoot[C]{\footnotesize\color{gray}-- Sổ tay kiến thức: Bài 08. Luyện tập chung --}
\renewcommand{\headrulewidth}{0.4pt}
\renewcommand{\footrulewidth}{0pt}

\setlength{\parindent}{0pt}
\setlength{\parskip}{0pt}

\begin{document}

% --- BANNER BÀI HỌC ---
\begin{tcolorbox}[
  enhanced,
  colback=teal!8!white,
  colframe=teal!70!black,
  arc=3mm, boxrule=1.2pt,
  left=3mm, right=3mm, top=2mm, bottom=2mm,
  halign=center
]
  {\LARGE\bfseries\color{teal!90!black}\faBookmark\ BÀI 08: LUYỆN TẬP CHUNG (ÔN TẬP BỔ SUNG)}\\[1mm]
  {\itshape\footnotesize\color{black!75} Tổng hợp kĩ năng cộng trừ 1000, tìm x, gấp và giảm số lần trên 1 trang A4}
\end{tcolorbox}

\vspace{-2mm}

\noindent
\begin{minipage}[t]{0.49\textwidth}\vspace{0pt}
\begin{tcolorbox}[
  enhanced,
  colback=blue!4!white,
  colframe=blue!75!black,
  arc=2mm, boxrule=0.8pt,
  left=2mm, right=2mm, top=1.8mm, bottom=1.8mm,
  title={\footnotesize\bfseries\color{white}\faCalculator\ 1. Tính toán \& Thứ tự phép tính}
]
{\footnotesize
\begin{itemize}\setlength{\itemsep}{1pt}\setlength{\leftskip}{-2.5mm}
  \item \textbf{Biểu thức chỉ có $(+, -)$ hoặc $(\times, :)$:} Thực hiện lần lượt từ \textbf{trái sang phải}.
  \item \textbf{Biểu thức có ngoặc tròn $( )$:} Thực hiện \textbf{trong ngoặc trước}, ngoài ngoặc sau.
  \item \textbf{Tính nhẩm với số tròn chục, tròn trăm:}
  $$300 + 400 = 700;\quad 800 - 500 = 300$$
\end{itemize}
}
\end{tcolorbox}
\end{minipage}
\hfill
\begin{minipage}[t]{0.49\textwidth}\vspace{0pt}
\begin{tcolorbox}[
  enhanced,
  colback=teal!4!white,
  colframe=teal!75!black,
  arc=2mm, boxrule=0.8pt,
  left=2mm, right=2mm, top=1.8mm, bottom=1.8mm,
  title={\footnotesize\bfseries\color{white}\faCheckSquare\ 2. Cẩm nang giải toán có lời văn}
]
{\footnotesize
\begin{itemize}\setlength{\itemsep}{1pt}\setlength{\leftskip}{-2.5mm}
  \item \textbf{Bước 1:} Đọc kĩ đề bài, xác định cái đã cho và cái cần tìm.
  \item \textbf{Bước 2:} Vẽ sơ đồ đoạn thẳng minh họa:
  \begin{itemize}\setlength{\itemsep}{0.5pt}\setlength{\leftskip}{-2.5mm}
    \item Tìm nhiều hơn / ít hơn $\to$ Vẽ đoạn dài hơn / ngắn hơn.
    \item Tìm gấp / giảm $\to$ Vẽ số phần bằng nhau.
  \end{itemize}
  \item \textbf{Bước 3:} Viết câu lời giải, phép tính và đáp số rõ ràng kèm tên đơn vị.
\end{itemize}
}
\end{tcolorbox}
\end{minipage}

\vspace{-2.5mm}
% --- KHỐI 3: SƠ ĐỒ TƯ DUY TIKZ ---
\noindent
\begin{tcolorbox}[
  enhanced,
  colback=orange!4!white,
  colframe=orange!80!black,
  arc=2mm, boxrule=0.8pt,
  left=2mm, right=2mm, top=1.5mm, bottom=1.5mm,
  title={\footnotesize\bfseries\color{white}\faSitemap\ SƠ ĐỒ TƯ DUY: HỆ THỐNG KIẾN THỨC CỐT LÕI CHỦ ĐỀ 1}
]
\centering
\begin{tikzpicture}[
  box/.style={draw, rounded corners=2pt, align=center, font=\scriptsize\bfseries, inner sep=2.5pt, minimum height=0.75cm},
  arr/.style={-{Stealth[length=2.5mm]}, thick}
]
  \node[box, fill=purple!15, draw=purple!80!black, text width=2.6cm] (c1) at (0, 0) {SỐ HỌC 1000\\Cấu tạo số, so sánh\\Cộng, trừ có nhớ};
  \node[box, fill=blue!12, draw=blue!80!black, text width=2.6cm] (c2) at (3.5, 0) {TÌM THÀNH PHẦN\\Số hạng = Tổng - SH\\SBT = Hiệu + ST};
  \node[box, fill=teal!12, draw=teal!80!black, text width=2.6cm] (c3) at (7.0, 0) {BẢNG 2 VÀ 5\\Nhân 2, 5 (Đếm thêm)\\Chia 2, 5 (Mối quan hệ)};
  \node[box, fill=orange!15, draw=orange!80!black, text width=2.6cm] (c4) at (10.5, 0) {GẤP / GIẢM\\Gấp $\to$ Nhân $(\times)$\\Giảm $\to$ Chia $(:)$};

  \draw[arr, purple!80!black] (c1) -- (c2);
  \draw[arr, blue!80!black] (c2) -- (c3);
  \draw[arr, teal!80!black] (c3) -- (c4);
\end{tikzpicture}
\end{tcolorbox}

\vspace{-2.5mm}
% --- KHỐI 4: VÍ DỤ & MẸO TRÁNH BẪY ---
\noindent
\begin{minipage}[t]{0.49\textwidth}\vspace{0pt}
\begin{tcolorbox}[
  enhanced,
  colback=green!4!white,
  colframe=green!60!black,
  arc=2mm, boxrule=0.8pt,
  left=2mm, right=2mm, top=1.8mm, bottom=1.8mm,
  title={\footnotesize\bfseries\color{white}\faCheckCircle\ Ví dụ thực hành chuẩn mực}
]
{\footnotesize
\textbf{1. Tính giá trị biểu thức:}\\
$150 + 200 - 70 = 350 - 70 = \mathbf{280}$.\\
$5 \times 8 - 15 = 40 - 15 = \mathbf{25}$.\\
\textbf{2. Tìm $x$:}\\
$x + 180 = 420 \implies x = 420 - 180 = \mathbf{240}$.\\
\textbf{3. Bài toán thực tế:}\\
Đoạn thẳng $AB = 8\text{ cm}$, đoạn $CD$ gấp $3$ lần đoạn $AB$. Đoạn $CD$ dài: $8 \times 3 = \mathbf{24}\text{ cm}$.
}
\end{tcolorbox}
\end{minipage}
\hfill
\begin{minipage}[t]{0.49\textwidth}\vspace{0pt}
\begin{tcolorbox}[
  enhanced,
  colback=red!4!white,
  colframe=red!70!black,
  arc=2mm, boxrule=0.8pt,
  left=2mm, right=2mm, top=1.8mm, bottom=1.8mm,
  title={\footnotesize\bfseries\color{white}\faExclamationTriangle\ Mẹo tránh bẫy trắc nghiệm}
]
{\footnotesize
\begin{itemize}\setlength{\itemsep}{1.5pt}\setlength{\leftskip}{-2.5mm}
  \item \textbf{Bẫy thứ tự nhân chia trước, cộng trừ sau:}\\
  Biểu thức $20 + 5 \times 2$: Phải tính $5 \times 2 = 10$ trước, rồi lấy $20 + 10 = 30$. Không được lấy $20 + 5 = 25$ rồi nhân $2$!
  \item \textbf{Bẫy quên đóng mở ngoặc đơn vị:}\\
  Trong câu trả lời, đơn vị đo phải đặt trong dấu ngoặc đơn, ví dụ: $24\text{ (cm)}$.
\end{itemize}
}
\end{tcolorbox}
\end{minipage}


\end{document}
